\documentclass{article}

% Same layout and visual style as report/template/template.tex and report no. 9.
\usepackage{xcolor}
\usepackage{tikz}
\usetikzlibrary{calc,positioning}
\usepackage{setspace}
\usepackage{geometry}
\usepackage{graphicx}
\usepackage{booktabs}
\usepackage{xepersian}

\linespread{1.0}
\geometry{margin=2.5cm}
\setlength{\parindent}{0.7cm}
\setlength{\textfloatsep}{10pt plus 2pt minus 2pt}

\definecolor{main_blue}{RGB}{0,70,140}
\definecolor{light_blue}{RGB}{0,120,200}
\definecolor{light_gray}{RGB}{112,146,190}

\newcommand{\reportnumber}[1]{\def\ReportNumber{#1}}
\newcommand{\reportdate}[1]{\def\ReportDate{#1}}
\newcommand{\reportstartdate}[1]{\def\ReportStartDate{#1}}
\newcommand{\reportenddate}[1]{\def\ReportEndDate{#1}}
\newcommand{\studentname}[1]{\def\StudentName{#1}}
\newcommand{\supervisorname}[1]{\def\SupervisorName{#1}}

\reportnumber{۱۰}
\reportdate{۱۴۰۵/۰۶/۲۸}
\reportstartdate{۱۴۰۵/۰۶/۱۸}
\reportenddate{۱۴۰۵/۰۶/۲۷}
\studentname{مهرداد شیخ عباسی}
\supervisorname{دکتر صاحب‌الزمانی}

\settextfont[
Scale=1.3,
Path=fonts/,
Extension=.TTF,
BoldFont=B-NAZANIN-Bold
]{B-NAZANIN}

\setlatintextfont[
Scale=1.1,
Path=fonts/,
Extension=.TTF,
BoldFont=timesbd,
ItalicFont=timesi
]{Times}

\renewcommand{\labelitemi}{\textcolor{main_blue}{\large$\bullet$}}

\newcommand{\mysection}[2]{
\par\vspace{0.7cm}
\noindent
\begingroup
\def\mysectionNumW{1.1cm}
\def\mysectionH{1.2cm}
\def\mysectionR{3pt}
\def\mysectionGap{0pt}
\begin{tikzpicture}
    \coordinate (rightedge) at (\dimexpr\textwidth-9pt\relax,0);
    \node[rectangle, fill=light_blue, text=white, anchor=east,
    minimum width=\mysectionNumW, minimum height=\mysectionH,
    rounded corners=\mysectionR, align=center, inner sep=2pt, outer sep=0pt,
    font=\Large\bfseries] (num) at (rightedge) {#1};
    \node[rectangle, fill=main_blue, text=white, anchor=east,
    minimum height=\mysectionH, rounded corners=\mysectionR,
    align=right, inner sep=6pt, outer sep=0pt,
    font=\Large\bfseries] (title) at ([xshift=-\mysectionGap]num.west) {#2};
    \draw[light_gray, line width=3pt, line cap=round]
    ($(0,0 |- num.north)+(0.5\pgflinewidth,0)$) --
    ($(rightedge |- num.north)+(-0.5\pgflinewidth,0)$);
\end{tikzpicture}
\endgroup
\par\vspace{0.35cm}
}

\newcommand{\reportheader}{
\vspace{0.5cm}
\noindent\makebox[\linewidth]{
{\Large گزارش هفتگی شماره \ReportNumber}
\hfill
{\Large \ReportDate}
}
\vspace{0.4cm}
\noindent{\color{light_gray}\rule{\linewidth}{0.5mm}}
\vspace{0.6cm}
\hspace{0.7cm}
\begin{minipage}{0.93\linewidth}
\textcolor{main_blue}{\textbf{تهیه‌کننده:}}~\StudentName
\hfill
\textcolor{main_blue}{\textbf{استاد راهنما:}}~\SupervisorName
\vspace{0.3cm}

\textcolor{main_blue}{\textbf{بازه زمانی:}}~\ReportStartDate~الی~\ReportEndDate
\end{minipage}
\vspace{0.4cm}
}

\begin{document}

\begin{figure}[h]
\centering
\includegraphics[width=4cm]{AUT.png}
\end{figure}
\reportheader

\noindent
با سلام و احترام،\\
هدف این مرحله، سنجش امکان اجرای مدل روی بردی با محیط \lr{PYNQ} بود. نخست \lr{Zynq-7020}، قطعه به‌کاررفته در برد \lr{PYNQ-Z2}، و سپس \lr{Zynq-7010} با منابع کمتر برای آزمون انتخاب شدند. این بررسی‌ها هنوز به انتخاب برد نهایی منجر نشده‌اند. در این بازه آزمون دقت جدیدی روی داده‌های آزمون انجام نشد و اعداد دقت گزارش پیشین تغییری نکرده‌اند.

\mysection{۱}{\rl{آماده‌سازی مدل و آزمون اولیه روی ۷۰۲۰}}

\noindent\textbf{مسیر داده.}
توابع بارگذاری و پیش‌پردازش داده و ساخت بارگذار\LTRfootnote{\lr{DataLoader}} بازبینی شدند. کنترل ورودی‌های نامعتبر و آزمون‌های مربوط به تقسیم در سطح فایل و پنجره‌سازی\LTRfootnote{\lr{Windowing}} تکمیل شد. تقسیم داده و تعریف برچسب‌ها تغییر نکرد. در اجرای فعلی، \lr{24/24} آزمون مرتبط گذر کرد و راهنمای مرور کد تا بخش مدل به‌روز شد.

\medskip
\noindent\textbf{آماده‌سازی هسته و سنتز\LTRfootnote{\lr{Synthesis}} اولیه.}
در هسته \lr{F\_1layer @ int10}، بافر\LTRfootnote{\lr{Buffer}} فعال‌سازی\LTRfootnote{\lr{Activation}} به ۱۰ بیت و انباره\LTRfootnote{\lr{Accumulator}} ضرب‌وجمع به ۳۲ بیت محدود شد. برخی بافرهای غیرهم‌زمان نیز دوباره استفاده شدند. برآورد \lr{Vitis HLS 2020.2} برای نسخه ردیفی با ضرایب صحیح \lr{Q40} روی قطعه \lr{XC7Z020CLG400-1} چنین است:

\begin{center}
\begin{latin}
\begin{tabular}{lrrrr}
\toprule
 & BRAM\_18K & LUT & FF & DSP \\
\midrule
Core estimate & 104 & 24,738 & 20,610 & 40 \\
Device capacity & 280 & 53,200 & 106,400 & 220 \\
\bottomrule
\end{tabular}
\end{latin}
\end{center}

\noindent
در این برآورد، مصرف هر چهار منبع هسته زیر ظرفیت قطعه است؛ اما منابع \lr{DMA} و دیگر اجزای سامانه، پیاده‌سازی نهایی و زمان‌بندی روی برد در آن نیامده‌اند. بنابراین از این جدول نمی‌توان جا شدن سامانه کامل روی ۷۰۲۰ را نتیجه گرفت.

\newpage
\noindent\textbf{کاهش حافظه میانی مدل.}
برای آزمودن قطعه کوچک‌تر، پردازش توکن‌ها\LTRfootnote{\lr{Tokens}} ردیفی شد. کلید\LTRfootnote{\lr{Key}} و مقدار\LTRfootnote{\lr{Value}} برای دنباله ۱۰۲۴ نمونه‌ای حفظ می‌شوند، اما پرس‌وجو\LTRfootnote{\lr{Query}} و مراحل پس از توجه\LTRfootnote{\lr{Attention}} هر بار برای یک ردیف محاسبه می‌شوند؛ بنابراین هر پرس‌وجو همچنان به کل دنباله دسترسی دارد. ضرایب باز‌مقیاس‌دهی\LTRfootnote{\lr{Requantization}} نیز به صورت صحیح \lr{Q40} درآمدند. شبیه‌سازی \lr{C} با طول اصلی دنباله، خروجی هر ۱۶ بردار مرجع\LTRfootnote{\lr{Golden vector}} را دقیقاً بازتولید کرد (\lr{16/16}).

برای اینکه برآورد منابع به همین نسخه نسبت داده شود، کد و بردارهای آزمون جداگانه ثبت شدند. اجرای نخست سنتز روی ۷۰۲۰ پس از حدود ۹۱ دقیقه بدون گزارش متوقف شد؛ اجرای دوباره همان نسخه کامل شد و اعداد آن در جدول صفحه پیش آمده‌اند. تأخیر تخمینی هر پنجره \lr{74,246,781} چرخه، معادل حدود \lr{0.742 s} با ساعت هدف \lr{10 ns} است. این عدد هنوز با هم‌شبیه‌سازی و آزمون روی برد تأیید نشده و برای ادعای بلادرنگ بودن کافی نیست.

\mysection{۲}{\rl{آزمون ۷۰۱۰ و نتیجه منابع}}

\noindent
سنتز نسخه ردیفی روی قطعه \lr{XC7Z010CLG400-1} در اجرای نخست بسیار طول کشید و گزارش نداد. خاموش‌کردن پایپ‌لاین\LTRfootnote{\lr{Pipeline}} حلقه‌ها نیز مشکل را رفع نکرد. پس از جداکردن محاسبات بعد از توجه در یک تابع، سنتز طول اصلی ۱۰۲۴ پس از حدود \lr{81.5} دقیقه کامل شد. شبیه‌سازی همین نسخه نیز هر چهار خروجی ۱۶ بردار مرجع را دقیقاً حفظ کرد.

جدول زیر تخمین منابع \emph{هسته کامل با طول اصلی} و ظرفیت قطعه ۷۰۱۰ را نشان می‌دهد:

\begin{center}
\begin{latin}
\begin{tabular}{lrrrr}
\toprule
 & BRAM\_18K & LUT & FF & DSP \\
\midrule
SEQ=1024 estimate & 104 & 24,518 & 21,072 & 40 \\
Z-7010 capacity & 120 & 17,600 & 35,200 & 80 \\
\bottomrule
\end{tabular}
\end{latin}
\end{center}

\noindent
مصرف \lr{LUT} معادل \lr{139\%} ظرفیت ۷۰۱۰ است؛ در نتیجه این نسخه فعلاً روی آن جا نمی‌شود. سه منبع دیگر در تخمین \lr{HLS} زیر ظرفیت‌اند. دوره ساعت تخمینی \lr{7.297 ns} در برابر هدف \lr{10 ns} است، اما تأخیر\LTRfootnote{\lr{Latency}} هر پنجره در گزارش نامشخص مانده و هنوز نمی‌توان درباره سرعت استنتاج قضاوت کرد.

برای بررسی راه کاهش منابع، ضرایب \lr{Q32} نیز آزمایش شدند. در شبیه‌سازی، ۱۶ بردار مرجع همان خروجی را دادند. در سنتز کوچک‌شده با طول ۶۴، مصرف \lr{LUT} از \lr{21,943} به \lr{20,991} رسید؛ این نتیجه را نمی‌توان به نسخه با طول ۱۰۲۴ تعمیم داد. نسخه \lr{Q32} هنوز روی داده آزمون کامل ارزیابی و برای طول اصلی سنتز نشده است.

گام بعد، کاهش مصرف \lr{LUT} در طول اصلی و تعیین تأخیر هر پنجره است. پس از آن باید هم‌شبیه‌سازی\LTRfootnote{\lr{Co-simulation}} \lr{C/RTL}، ساخت \lr{bitstream} و آزمون روی برد انجام شود. این مراحل هنوز تکمیل نشده‌اند.

\end{document}
